\clearpage
\section{Özet}
\label{CH:Ozet}

Bu projede, işitme engelli bireylerle işaret dili bilmeyen kişiler arasındaki iletişimi
kolaylaştırmak amacıyla, akıllı telefon kamerasından alınan Türk İşaret Dili (TİD)
kelimelerini tanıyıp yazıya çeviren bir Android uygulaması geliştirilmiştir. Sistem video
karesini değil, telefonda çıkarılan iskelet koordinatlarını kullanmaktadır. Bu sayede
internet bağlantısına ihtiyaç duyulmamakta ve kamera görüntüsü cihazdan dışarı
çıkmamaktadır.

Eğitim ve değerlendirmede AUTSL veri setinin işaretçiden bağımsız resmî bölmeleri
kullanılmıştır (226 işaret, 43 işaretçi). Her karede iki elin 21'er noktası ve üst gövdeye
ait 7 nokta Google MediaPipe Tasks ile çıkarılmış, koordinatlar omuz ortasına ötelenip
omuz genişliğine bölünmüş ve diziler 64 kareye örneklenmiştir. Aynı veri ve eğitim
düzeniyle iki model karşılaştırılmıştır: iki katmanlı, iki yönlü GRU (BiGRU, 2,11 milyon
parametre) ve 1B evrişim ile iki katmanlı Transformer kodlayıcısından oluşan model
(1B-CNN + Transformer, 1,47 milyon parametre). Üç farklı tohumla yapılan 60 epoch'luk
eğitimler sonunda 3.741 videoluk test kümesinde BiGRU yüzde $87{,}98 \pm 0{,}23$,
1B-CNN + Transformer yüzde $89{,}42 \pm 0{,}52$ Top-1 doğruluğa ulaşmış ve tez
önerisindeki yüzde 85 hedefini aşmıştır.

Seçilen 1B-CNN + Transformer modeli Keras üzerinden LiteRT biçimine aktarılmış ve
dinamik aralık int8 nicemleme uygulanmıştır. Model dosyası 1,61 MB'tır; test
doğruluğu nicemleme sonrası yüzde 89,82'dir ve PyTorch modeliyle tahminlerin yüzde
99,57'si aynıdır. Mobil uygulama Flutter ve Kotlin ile geliştirilmiştir: kamera,
MediaPipe ve model çıkarımı Kotlin tarafında çalışmakta, arayüz canlı tanıma, metin ve
işaret listesi olmak üzere üç ekrandan oluşmaktadır. İşaret listesinde her kelimenin
eğitim verisinden alınmış örnek hareketi iskelet animasyonu olarak gösterilmektedir.

Kaydedilmiş veri üzerinde hedefler sağlanırken iki konu açık kalmıştır. Uçtan uca hız,
iki MediaPipe modelinin kare başına çalışması nedeniyle 15 FPS hedefinin altındadır;
canlı kamera kullanımında ise tahmin güvenleri kayıtlı veriye göre düşmektedir. Bu iki
konuyu ölçmek ve ayrıştırmak için uygulamaya bir gönüllü testi ekranı ve kayıtları
bilgisayarda inceleyen bir analiz aracı eklenmiştir.

\vspace{3mm}
\noindent\textbf{Anahtar kelimeler:} Türk İşaret Dili, işaret dili tanıma, MediaPipe,
iskelet tabanlı tanıma, Transformer, çevrimdışı mobil uygulama, LiteRT.

\clearpage
